% !TeX root = ../../main.tex

Die in Abschnitt~\ref{sec:errorclasses} beschriebenen Fehlerklassen können dazu führen, dass Transaktionen in
PostgreSQL nicht erfolgreich abgeschlossen werden können.
Besonders bei parallelen Zugriffen mehrerer Prozesse oder Anwendungen auf gemeinsame Datenbankressourcen können
Konflikte entstehen, die eine erfolgreiche Durchführung einzelner Transaktionen verhindern.
Diese Situationen treten beispielsweise auf, wenn mehrere Transaktionen gleichzeitig versuchen einen Datensatz zu
verändern.
Das Datenbanksystem bricht daraufhin eine oder mehrere Transaktionen ab, um die Konsistenz der Daten sicherzustellen.

Für Anwendungen bedeutet ein solcher Abbruch, dass die fehlgeschlagene Transaktion erneut ausgeführt werden muss.
Während dies bei einfachen Anwendungen, wie bei einzelnen vom Nutzer ausgeführten Datenbankoperationen, durch einen
manuellen Neustart der Transaktion erfolgen kann, stellt dies in komplexeren Systemen eine Herausforderung dar.
Moderne Anwendungen bestehen häufig aus mehreren abhängigen Komponenten, die gleichzeitig auf gemeinsame Daten
zugreifen.
Tritt in dieser Kette an Komponenten an einer Stelle ein Fehler auf, kann dieser Fehler dafür sorgen, dass weitere
Prozesse ebenfalls fehlschlagen oder nicht mehr korrekt ausgeführt werden.
Besonders kritisch ist dies bei Systemen, bei denen Transaktionen automatisch und ohne Nutzereingriffe verarbeitet
werden.

Um solche Fehlerfälle automatisiert zu behandeln, werden Retry-Mechanismen eingesetzt.
Die automatisierte Behandlung fehlgeschlagener Transaktionen trägt dazu bei, die Verfügbarkeit von Anwendungen zu
erhöhen und manuelle Eingriffe zu reduzieren.
Gleichzeitig darf die Fehlerbehandlung jedoch nicht zu einer unverhältnismäßigen Belastung des Datenbanksystems
führen.
Der folgende Abschnitt führt die Risiken und Grenzen von Retry auf, auf die bei einer Implementierung geachtet
werden muss.
